\section{Introduction}
Frozen vision-language models (VLMs) provide transferable image--text representations for classification and retrieval \citep{radford2021clip}. Their zero-shot interface, however, does not remove biases inherited from web-scale pretraining. Errors can concentrate on demographic groups \citep{agarwal2021evaluating,hall2023gender}, but also on backgrounds, acquisition conditions, pose, or other spurious factors. Standard average accuracy can conceal these failures.

Most post-hoc VLM debiasing methods assume that the relevant attribute is known. They encode named concepts in prompts \citep{chuang2023debiasing}, use attribute-labeled reference images \citep{gerych2024bendvlm}, or estimate attribute-predictive subspaces \citep{zhao2026subspace}. This is appropriate when the attribute is available and is the actual source of harm. In many deployments, however, group labels are unavailable, restricted, incomplete, or simply miss the factor driving errors. What remains observable in deployment is the model’s pattern of errors and low-confidence predictions.

We therefore ask: \emph{can low classification margins—where the correct-class score is close to that of the strongest competing class—and prediction errors reveal a representation direction whose removal reduces hidden-group failure?} This question differs from discovering a demographic attribute. A failure-predictive direction may encode background, image quality, an intersection of factors, or useful class evidence. Consequently, editing it may improve worst-group reliability, do nothing, or make performance less balanced. A credible method must expose this distinction rather than equating any failure direction with a fairness direction.

We call a low-dimensional subspace a \emph{latent risk subspace} when it is predictive of within-class calibration failures, weakly coupled to class semantics, and suppressible without destroying the task score. This is an operational representation of error concentration, not a recovered gender, race, age, or other sensitive attribute. If it overlaps the factor defining a worst-performing group, suppressing it can indirectly improve worst-group performance; if it does not, editing can be harmful.

We introduce \textsc{LatentRisk-VLM}, a post-hoc framework for frozen VLM embeddings. Given a task-labeled calibration set, we derive risk labels from zero-shot margins and remove class prototypes and text-aligned evidence. We estimate a class-conditional risk direction either by contrasting high- and low-risk residual means or by fitting iterative linear probes inspired by INLP \citep{ravfogel2020inlp}. Test labels are not used: each candidate class supplies its own prototype and subspace, producing candidate-wise corrected scores. All selection in the primary protocol uses task labels only; sensitive attributes are read after prediction solely for group metrics.

The completed experiments give a deliberately mixed answer. On Waterbirds, rank-1 editing improves WGA by 26.9 points without place labels, although it loses 5.3 points of average accuracy. Rank-4 editing with group-free worst-class validation selection reaches 63.1 WGA on Waterbirds and 72.8 WGA on CelebA Blond Hair. However, other group-free selectors over the same candidate pool perform much worse, and attribute-labeled validation still raises rank-4 CelebA WGA to 82.2. A post-hoc diagnostic explains the rank-1 contrast: Waterbirds risk scores are enriched for minority-place groups, while CelebA risk directions have little alignment with the Male attribute direction. A frozen-SigLIP extension preserves the cautionary part of the story: task-accuracy selection fails on both datasets, while rank-4 worst-class selection improves WGA relative to a weak frozen SigLIP baseline but often pays a large average-accuracy cost. On a MetaShift-style cat/dog indoor/outdoor shift, a preregistered diagnostic predicts uncertainty and the final held-out gain is positive but small (+5.6 WGA points). Discovering a risk direction and selecting a group-robust operating point are therefore separate problems; CnC-Frozen with attribute-labeled validation reaches 85.3. A FairFace/Flickr30K retrieval audit likewise finds a modest skew reduction accompanied by a 4.0-point Recall@10 loss.

Our contributions are fivefold. First, we formulate latent-risk editing in frozen VLMs using task labels and failure signals, but no group or sensitive-attribute annotations, to estimate class-conditional risk directions. At test time, because the true class is unknown, the method edits and scores each candidate class separately rather than selecting an editor using the ground-truth label. Second, we develop an attribute-label-free, post-hoc low-rank editing procedure that leaves both VLM encoders unchanged and runs on cached features. Third, we add an attribute-label budget protocol that exposes the zero-annotation operating point and separates training supervision from model-selection supervision. Fourth, we compare against supervised concept-erasure baselines such as INLP and LEACE to test when failure-derived directions substitute for a named attribute direction and when they do not. Fifth, under a shared frozen-backbone setting, we compare the additional supervision, optimization, and adaptation cost required after feature extraction. The contrasting results are therefore central to our conclusion: failures provide useful weak supervision when they align with hidden groups, but they do not consistently identify those groups.
